\src{186}
Now consider the pairings defined by composition of $\operatorname{Ext}$:
\begin{equation}\tag{6.4}
  H^{p}(X,F)\times
  \operatorname{Ext}_{\Osh_X}^{n-p}(X;F,\omega_X^{n})
  \longrightarrow H^{n}(X,\omega_X^{n}),
\end{equation}
and compose them with the homomorphism $\gamma$ of (6.3). This gives functorial
homomorphisms, compatible with the boundary operators,
\begin{equation}\tag{6.5}
  \operatorname{Ext}_{\Osh_X}^{n-p}(X;F,\omega_X^{n})
  \longrightarrow H^{p}(X,F)'
\end{equation}
(generalizing (5.2)). One checks that for $p=n$ this gives the isomorphism
(6.1). With this understood, we have the following result.

\result{THEOREM}{3 bis} The following four conditions on $X$ are equivalent for
a given integer $k\geq0$:
\begin{enumerate}
  \item[\textit{i.}] The functorial homomorphism (6.5) is an isomorphism for
    $n-k\leq p\leq n$.
  \item[\textit{ii.}] One has $H^{p}(X,\Osh_X(-m))=0$ for $m$ large and
    $n-k\leq p<n$.
  \item[\textit{iii.}] The functor $H^{p}(X,F)$ on the category of coherent
    algebraic sheaves on $X$ is coeffaceable for $n-k\leq p<n$.
  \item[\textit{iv.}] One has
    $E^{i}(\Osh_X)=\Extsh_{\Osh_P}^{r-n+i}(\Osh_X,\Omega_P^{r})=0$
    for $0<i\leq k$.
\end{enumerate}

\noindent\textbf{PROOF.} --- \textit{i.}$\Rightarrow$\textit{ii.} follows
from Lemma 4, \textit{ii.}$\Rightarrow$\textit{iii.} from Lemma 3, and
\textit{iii.}$\Rightarrow$\textit{i.} by a well-known standard argument, since
the two sides of (6.5) are then coeffaceable functors for $n-k\leq p<n$ (the
first by Lemma 4). Finally, let us prove
\textit{ii.}$\Leftrightarrow$\textit{iv.} This follows from the corollary to the
next proposition.

\result{PROPOSITION}{6} Let $F$ be a coherent algebraic sheaf on $X$, and let
$i$ be an integer. Then, for $m$ sufficiently large, there is an isomorphism
\begin{equation}\tag{6.6}
  H^{i}(X,F(-m))'
  \simeq H^{0}\bigl(X,E^{n-i}(F)(m)\bigr),
\end{equation}
where we set (compare no. 3, Proposition 5)
\begin{equation}\tag{6.7}
  E^{j}(F)=\Extsh_{\Osh_P}^{r-n+j}(F,\Omega_P^{r}).
\end{equation}
Indeed, by the duality theorem on $P$, the first side of (6.6) is isomorphic to
$\operatorname{Ext}_{\Osh_P}^{r-i}(P;F(-m),\Omega_P^{r})$, so (6.6) follows
from Lemma 5.

